%% ============================================================================================================
%% Data and Study Area
%% ============================================================================================================
\section{Data and Study Area}
\label{sec:data}

% Study Area
% ————————————————————————————————————————————————————————————————————————————
\subsection{Study Area}

This study focuses on river basins across Eurasia, located approximately between latitudes \todo{latitudes} and longitudes \todo{longitudes}. The selected region encompasses a wide range of static basin characteristics and climatic conditions, making it suitable for analyzing hydrological equifinality.

\todo[inline]{Add a map of the river basins}

% Hydrometeorological Data
% ————————————————————————————————————————————————————————————————————————————
\subsection{Hydrometeorological Data}

The dataset includes daily time series from 2008 to 2022 for the following features:
precipitation (\texttt{prcp}), river water level (\texttt{lvl\_sm}), minimum temperature (\texttt{t\_min}), maximum temperature (\texttt{t\_max}), mean temperature (\texttt{t\_mean} \todo{check the name of t\_mean}), and river discharge (\texttt{q\_mm\_day}).

\todo[inline]{Add an example time series plot.}

% Static Basin Attributes
% ————————————————————————————————————————————————————————————————————————————
\subsection{Static Basin Attributes}

Static (time-invariant) characteristics of the river basins were derived from the HydroSHEDS database \todo{Add cite to https://www.hydrosheds.org/}. A total of 12 features\ref{table:static-attributes} were selected based on two main criteria: (i) low pairwise correlation and (ii) high interpretability for ability to explain the models.

In total, data were available for 2,201 river basins. For model training and subsequent analysis, we selected a subset of 1,073 basins with the fewest missing values to ensure data quality and consistency across samples.

\begin{table}[ht]
    \centering
    \begin{tabular}{lll}
        \toprule
        \textbf{ID}  & \textbf{Attribute}              & \textbf{Column Name}  \\
        \midrule
        \textbf{L07} & Forest Cover Extent             & \texttt{for\_pc\_sse} \\
        \textbf{L08} & Cropland Extent                 & \texttt{crp\_pc\_sse} \\
        \textbf{L09} & Pasture Extent                  & \texttt{pst\_pc\_sse} \\
        \textbf{L10} & Irrigated Area Extent           & \texttt{ire\_pc\_sse} \\
        \textbf{A03} & Urban Extent                    & \texttt{urb\_pc\_sse} \\
        \textbf{S07} & Karst Area Extent               & \texttt{kar\_pc\_sse} \\
        \textbf{P01} & Elevation                       & \texttt{ele\_mt\_sav} \\
        \textbf{P02} & Terrain Slope                   & \texttt{slp\_dg\_sav} \\
        \textbf{C06} & Aridity Index                   & \texttt{ari\_ix\_sav} \\
        \textbf{H10} & Groundwater Table Depth         & \texttt{gwt\_cm\_sav} \\
        \textbf{S05} & Soil Water Content (Annual Avg) & \texttt{swc\_pc\_syr} \\
        \textbf{S01} & Clay Fraction in Soil           & \texttt{cly\_pc\_sav} \\
        \textbf{A06} & Human Footprint Index           & \texttt{hft\_ix\_s09} \\
        \textbf{P04} & Latitude                        & \texttt{lat}          \\
        \textbf{P05} & Longitude                       & \texttt{lon}          \\
        \bottomrule
    \end{tabular}
    \caption{List of selected static features and their corresponding column names.}
    \label{table:static-attributes}
\end{table}

